% Historical two-model focused recovery and separate ChatGPT round-2 evaluation.
\begin{abstract}
LLM-generated JUnit suites can fail because their assertions or fixtures do not match the evaluated Java source. We examine a failure-informed Oracle-Aware Prompt Refinement (OAPR) workflow on 94 ClassEval-T classes and two reported model configurations. After excluding Gemini CE023, whose P0 prompt used a different target source, and correcting two ChatGPT dependency failures without changing their tests, 187 model--task pairs remain. The adjusted P0 baseline passed 169/187. One archived P1 regeneration recovered 15 of 18 selected failures (83.33\%; exact 95\% interval 58.58--96.42\%): 10/11 for ChatGPT and 5/7 for Gemini. Carrying forward P0 passes gives 184/187 (98.40\%); this is not a full-P1 result. A separate ChatGPT round-2 evaluation recorded three runs of each of three conditions on 11 selected tasks. P0 resampling passed 23/33 suites; checklist-only and checklist-plus-focus each passed 33/33. The two checklist conditions had identical recorded mutation outcomes (924/1,164 mutants killed each), while P0 mutation scoring included only its 23 baseline-valid suites. Incomplete generation provenance, repeated observations of the same tasks, and the absence of Gemini round-2 arms limit inference. These results describe observed executable compatibility and focused recovery, not an isolated causal effect or general defect-detection improvement.
\end{abstract}

\begin{IEEEkeywords}
large language models, unit test generation, test oracle, prompt refinement, JUnit, ClassEval-T
\end{IEEEkeywords}
